Magnetic microscopy allows scientists to measure tiny magnetic signals from individual mineral grains and small geological samples. These measurements can provide valuable information about the Earth's ancient magnetic field and the magnetic properties of rocks. As magnetic microscopes continue improving, they can now measure signals at higher resolutions and closer distances to the sample surface. However, these advances also reveal more complex magnetic patterns that are often too complicated to be accurately represented by traditional dipole-based methods. In this study, we developed a new approach that combines spherical harmonic models with an automatic estimate of the magnetic source location using Euler deconvolution. The method was tested using both computer-generated magnetic models and real magnetic microscopy data. The results show that the new approach more accurately reproduces complex magnetic signals, especially when measurements are acquired close to the sample surface where non-dipolar magnetic effects become important. This work provides a practical way to extract more information from modern magnetic microscopy measurements and may improve future studies of magnetic minerals, paleomagnetism, and rare geological materials.